% GENERATED FILE. Edit canonical lessons, then run npm run generate:books.
% canonical-source-hash: fnv1a64:a2673b506dd16453
% canonical-lessons: AR-C66-aruzz, AR-C66-fakiha, AR-C66-tuffaha, AR-C66-mawz, AR-C66-batatis

\chapter{Rice, Fruit, Apple, Bananas, Potatoes}
\label{ch:ar-rice-fruit-apple-bananas-potatoes}
\hlchaptermodality{66}

\begin{chapteropening}
\textbf{By the end of this chapter:} \emph{I can say rice, fruit, an apple, bananas and potatoes.}

Complete the last lesson of chapter 66: I can say rice, fruit, an apple, bananas and potatoes.
\end{chapteropening}

\section[\texorpdfstring{\ar{أرز} (aruzz)}{أرز (aruzz)}]{\ar{أرز} (aruzz) — rice}

\label{lesson:AR-C66-aruzz}



\begin{quote}
Before the new one: say the Arabic for to go down, then the Arabic for fish.
\end{quote}

\subsection*{You'll want to know: \ar{أرز}}
\textbf{\ar{أرز}} — \emph{aruzz} — \textquotedblleft{}rice\textquotedblright{}.

\begin{cousinweb}
\textbf{\ar{أرز}} (\emph{aruzz}) is an old loan. It is usually traced, probably by way of Aramaic, to Greek \emph{oryza}, \textquotedblleft{}rice\textquotedblright{} (EtymDB), a word the Greeks had from further east. Lane lists \emph{aruzz} under \emph{ʾ-r-z} as a well-known grain.
\end{cousinweb}

\begin{grammarlens}[title={What you've built}]
One more word for this chapter.
\end{grammarlens}

\subsection*{Your turn}
\begin{itemize}
\raggedright
  \item \emph{Say it:} \emph{aruzz}
  \item \emph{Say it:} \emph{aruzz}, once more
  \item \emph{Say it:} say \emph{aruzz}
  \item \emph{Recall:} say the Arabic for to go down, then the Arabic for fish, then say \emph{aruzz} again
\end{itemize}

\begin{tcolorbox}[breakable,colback=teal!4,colframe=teal!35!black,title={Before you move on}]
What does \ar{أرز} mean? (\textquotedblleft{}Rice\textquotedblright{}.) Say it once more.
\end{tcolorbox}
\section[\texorpdfstring{\ar{فاكهة} (fākiha)}{فاكهة (fākiha)}]{\ar{فاكهة} (fākiha) — fruit}

\label{lesson:AR-C66-fakiha}



\begin{quote}
Before the new one: say the Arabic for fish, then the Arabic for rice.
\end{quote}

\subsection*{You'll want to know: \ar{فاكهة}}
\textbf{\ar{فاكهة}} — \emph{fākiha} — \textquotedblleft{}fruit\textquotedblright{}.

\begin{cousinweb}
\textbf{\ar{فاكهة}} (\emph{fākiha}) is built on the native Arabic root \emph{f-k-h}, \textquotedblleft{}fruit; to enjoy\textquotedblright{} (Lane). Lane explains \emph{fākiha} as fruit of any kind, a thing eaten for enjoyment.
\end{cousinweb}

\begin{grammarlens}[title={What you've built}]
One more word for this chapter.
\end{grammarlens}

\subsection*{Your turn}
\begin{itemize}
\raggedright
  \item \emph{Say it:} \emph{fākiha}
  \item \emph{Say it:} \emph{fākiha}, once more
  \item \emph{Say it:} say \emph{fākiha}
  \item \emph{Recall:} say the Arabic for fish, then the Arabic for rice, then say \emph{fākiha} again
\end{itemize}

\begin{tcolorbox}[breakable,colback=teal!4,colframe=teal!35!black,title={Before you move on}]
What does \ar{فاكهة} mean? (\textquotedblleft{}Fruit\textquotedblright{}.) Say it once more.
\end{tcolorbox}
\section[\texorpdfstring{\ar{تفاحة} (tuffāḥa)}{تفاحة (tuffāḥa)}]{\ar{تفاحة} (tuffāḥa) — an apple}

\label{lesson:AR-C66-tuffaha}



\begin{quote}
Before the new one: say the Arabic for rice, then the Arabic for fruit.
\end{quote}

\subsection*{You'll want to know: \ar{تفاحة}}
\textbf{\ar{تفاحة}} — \emph{tuffāḥa} — \textquotedblleft{}an apple\textquotedblright{}.

\begin{cousinweb}
\textbf{\ar{تفاحة}} (\emph{tuffāḥa}) is built on the native Arabic root \emph{t-f-ḥ}, \textquotedblleft{}the apple\textquotedblright{} (Lane). Lane quotes a source calling it an Arabic word, not arabicized; Hebrew \emph{tappuaḥ}, \textquotedblleft{}apple\textquotedblright{}, is its cousin. \emph{Tuffāḥa} is one apple.
\end{cousinweb}

\begin{grammarlens}[title={What you've built}]
One more word for this chapter.
\end{grammarlens}

\subsection*{Your turn}
\begin{itemize}
\raggedright
  \item \emph{Say it:} \emph{tuffāḥa}
  \item \emph{Say it:} \emph{tuffāḥa}, once more
  \item \emph{Say it:} say \emph{tuffāḥa}
  \item \emph{Recall:} say the Arabic for rice, then the Arabic for fruit, then say \emph{tuffāḥa} again
\end{itemize}

\begin{tcolorbox}[breakable,colback=teal!4,colframe=teal!35!black,title={Before you move on}]
What does \ar{تفاحة} mean? (\textquotedblleft{}An apple\textquotedblright{}.) Say it once more.
\end{tcolorbox}
\section[\texorpdfstring{\ar{موز} (mawz)}{موز (mawz)}]{\ar{موز} (mawz) — bananas}

\label{lesson:AR-C66-mawz}



\begin{quote}
Before the new one: say the Arabic for fruit, then the Arabic for an apple.
\end{quote}

\subsection*{You'll want to know: \ar{موز}}
\textbf{\ar{موز}} — \emph{mawz} — \textquotedblleft{}bananas\textquotedblright{}.

\begin{cousinweb}
\textbf{\ar{موز}} (\emph{mawz}) is a loan. EtymDB traces it to Middle Persian \emph{mōz}, which is usually taken back to Sanskrit \emph{mocā}, \textquotedblleft{}banana\textquotedblright{}. Lane lists it under \emph{m-w-z} as a well-known tree and its fruit.
\end{cousinweb}

\begin{grammarlens}[title={What you've built}]
One more word for this chapter.
\end{grammarlens}

\subsection*{Your turn}
\begin{itemize}
\raggedright
  \item \emph{Say it:} \emph{mawz}
  \item \emph{Say it:} \emph{mawz}, once more
  \item \emph{Say it:} say \emph{mawz}
  \item \emph{Recall:} say the Arabic for fruit, then the Arabic for an apple, then say \emph{mawz} again
\end{itemize}

\begin{tcolorbox}[breakable,colback=teal!4,colframe=teal!35!black,title={Before you move on}]
What does \ar{موز} mean? (\textquotedblleft{}Bananas\textquotedblright{}.) Say it once more.
\end{tcolorbox}
\section[\texorpdfstring{\ar{بطاطس} (baṭāṭis)}{بطاطس (baṭāṭis)}]{\ar{بطاطس} (baṭāṭis) — potatoes}

\label{lesson:AR-C66-batatis}



\begin{quote}
Before the new one: say the Arabic for an apple, then the Arabic for bananas.
\end{quote}

\subsection*{You'll want to know: \ar{بطاطس}}
\textbf{\ar{بطاطس}} — \emph{baṭāṭis} — \textquotedblleft{}potatoes\textquotedblright{}.

\begin{cousinweb}
\textbf{\ar{بطاطس}} (\emph{baṭāṭis}) is a modern loan from a European name of the potato, usually traced to Spanish \emph{patata}, which goes back to Taíno \emph{batata}, \textquotedblleft{}sweet potato\textquotedblright{}. It has no Arabic root.
\end{cousinweb}

\begin{grammarlens}[title={What you've built}]
One more word for this chapter.
\end{grammarlens}

\subsection*{Your turn}
\begin{itemize}
\raggedright
  \item \emph{Say it:} \emph{baṭāṭis}
  \item \emph{Say it:} \emph{baṭāṭis}, once more
  \item \emph{Say it:} say \emph{baṭāṭis}
  \item \emph{Recall:} say the Arabic for an apple, then the Arabic for bananas, then say \emph{baṭāṭis} again
\end{itemize}

\begin{tcolorbox}[breakable,colback=teal!4,colframe=teal!35!black,title={Before you move on}]
What does \ar{بطاطس} mean? (\textquotedblleft{}Potatoes\textquotedblright{}.) Say it once more.
\end{tcolorbox}
